%%
%% Copyright (c) 2018-2020 Weitian LI <wt@liwt.net>
%% CC BY 4.0 License
%%
%% Created: 2018-04-11
%%

% Chinese version
\documentclass[zh]{resume}

% === 强制一页:紧缩间距 ===
\usepackage{geometry}
\geometry{margin=0.9cm}           % 默认 1.5cm → 0.9cm
\linespread{1.1}               % 中文默认 1.25 → 1.1
\titlespacing{\section}{0em}{0.5ex}{0.5ex}  % 节前后间距减半
\setlist[itemize,1]{label=\faAngleRight, nosep, leftmargin=2em, topsep=1pt}
\setlength{\medskipamount}{4pt}

% File information shown at the footer of the last page
\fileinfo{%
  % \faCopyright{} 2025--2026, 阿东玩AI \hspace{0.5em}
  % \creativecommons{by}{4.0} \hspace{0.5em}
  % \githublink{dok4everak47}{LLM-Resume-Template} \hspace{0.5em}
  % \faEdit{} \today
}

\name{劲舟}{梁}

\tagline{求职意向: Agent Harness 方向}

\photo[shape=circular, position=right]{1.5cm}{../assets/luongchin.png}

\profile{
  \mobile{18176608865}
  \email{dok4ever123@126.com}
  \iconlink{\faGithub}{https://github.com/dok4everak47}{dok4everak47}
}

\begin{document}
\begin{tcolorbox}[colback=accentcolor!12, colframe=accentcolor!20,
  arc=0pt, boxrule=0.3pt, left=6pt, right=6pt, top=3pt, bottom=3pt,
  sharp corners]
\makeheader
\end{tcolorbox}

\vspace{0.3em}

%======================================================================
\sectionTitle{个人总结}{\faUser}
%======================================================================
具有数学与计算科学背景的 Agent Harness 研究者。对 LLM Agent 的 Agent Loop、Tool Use、Memory、Context
Engineering 机制有深入研究，并通过 claude-code-from-scratch 项目从零复现了 Coding Agent 核心架构。熟悉
TypeScript 和 Python，能动手构建 Agent 工具链和评测分析系统。热衷于拆解 Agent 行为、分析失败模式、推动 Agent 系统
持续进化。
% \par\vspace{0.3ex}
% English: Fluent reading and writing — able to read technical papers and documentation,
% follow overseas AI communities and open-source projects.

%======================================================================
\sectionTitle{专业技能}{\faWrench}
%======================================================================
\begin{itemize}
    \item \textbf{编程语言}: TypeScript, Python, Node.js, Pandas, NumPy, scikit-learn
    \item \textbf{Agent 工程}: Agent Loop / Tool Use / Prompt Engineering / MCP 协议 / Memory 系统 / Context Engineering
    \item \textbf{大模型}: RLHF 数据构建与质量评估、模型输出分析、多模态 AIGC 评估
    \item \textbf{工具链}: Git, CLI 工具开发, 日志分析, Vercel Serverless 部署, API 集成 (Anthropic / OpenAI / DeepSeek)
\end{itemize}

%======================================================================
\sectionTitle{个人项目}{\faRocket}
%======================================================================
\begin{itemize}
            \item[\faCode] \textbf{Claude Code from Scratch — Coding Agent 核心架构复现与扩展}

    从零实现 Claude Code 核心的 Agent Harness 架构，涵盖 Agent Loop、工具系统、上下文压缩、记忆系统、多 Agent 架构等核心模块，总代码量约 4300 行。

    技术栈: TypeScript, Node.js, Anthropic / OpenAI API \\
    GitHub: \href{https://github.com/dok4everak47/claude-code-from-scratch}{\texttt{dok4everak47/claude-code-from-scratch}}
    \begin{itemize}
        \item 完整构建 Agent 执行循环与 13 个内置工具，支持流式并行执行、跨 Provider 消息排序安全网、分层上下文压缩、多类型记忆系统、MCP 协议集成
        \item 开发 Agent Tool System Demo（React + Vite + Tailwind CSS）：4 种预设场景逐帧展示 Agent 推理与安全网恢复；自由模式连接真实 LLM API，集成 Wikipedia 百科/实时汇率/词典定义/趣味笑话等公开数据工具，Tool 按阶段过滤，实时面板展示运行状态；对比模式并行 3 种 System Prompt 对比策略差异及量化评测指标；接入 Open-Meteo 天气 API 并通过 Vercel Proxy 保护 API Key
    \end{itemize}

    \vspace{1ex}
    {\color{accentcolor!15}\hrule}
    \vspace{1ex}

    \item[\faCode] \textbf{Agent Harness 行为分析与评测工具链}

    自建 Agent 对话日志分析系统，使用 TypeScript + Python（Pandas / NumPy / scikit-learn）处理对话日志，解决多源日志交错解析、噪声清洗等工程问题。

    技术栈: TypeScript, Python, Pandas, NumPy, scikit-learn
    \begin{itemize}
        \item 基于日志分析提出三层分级记忆模型与动态摘要缓存方案，对应 Harness 中的 Memory 和 Context Engineering 优化方向
        \item 对 Pi / Claude / ChatGPT / DeepSeek 进行对比测试，建立 Agent 行为对比评测框架，产出「情绪-风格解耦对齐」「阈值触发式风格切换」「跨对话情绪」3 项可量化洞察
        \item 分析脚本可扩展用于任意 Agent 对话日志的行为模式提取与量化分析
    \end{itemize}
\end{itemize}

%======================================================================
\sectionTitle{工作经历}{\faBriefcase}
%======================================================================
\begin{experiences}
  \experience%
    [2025.12]%
    {2026.05}%
    {广西一三数据科技有限公司 \quad AI训练师}%
    [\begin{itemize}
      \item 主导多模态 AIGC 数据评估体系的构建，制定覆盖 15+ 类场景的评估标准，将团队交付合格率提升至 98\%
      \item 负责模型输出的质量分析与人工蒸馏，累计产出高质量偏好数据 100 万+条，支撑模型 RLHF 训练与能力迭代
      \item 主导自动驾驶 3D 感知数据的标注策略制定，处理叠帧重影和移动物体轨迹连续性，确保 Agent 环境感知数据的质量一致性，累计处理 100 万+ 目标物，支撑自动驾驶模型环境理解精度提升
    \end{itemize}]
\end{experiences}

%======================================================================
\sectionTitle{教育背景}{\faGraduationCap}
%======================================================================
\begin{educations}
  \education%
    {2020.09}%
    [2023.06]%
    {桂林电子科技大学}%
    {数学与计算科学学院}%
    {数学}%
    {硕士}
  \separator{0.2ex}
  \education%
    {2016.09}%
    [2020.06]%
    {西安邮电大学}%
    {理学院}%
    {信息与计算科学}%
    {学士}
\end{educations}

\end{document}